\tsDef{(Series)}{
    A (formal) expression of the form
    \(
        a_0 + a_1 + a_2 + \cdots = \sum_{n=0}^{\infty} a_n
    \)
    is called a \emph{series}. The terms $a_n$ are called elements or summands.
    If the sequence of partial sums
    \(
        s_n := \sum_{k=0}^{n} a_k
    \)
    converges to a limit $L \in \mathbb{R}$, then the series is said to \emph{converge}, and we write
    \(
        \sum_{n=0}^{\infty} a_n = L.
    \)
    Otherwise, the series is called \emph{divergent}.
}
\\
\tsThe{(Necessary Condition for Convergence)}{
    If the series $\sum_{n=0}^{\infty} a_n$ converges, then necessarily
    \(
        \lim_{n\to\infty} a_n = 0.
    \)
}
\\
\tsLem{(Linearity of Series)}{
    Let $\sum_{n=0}^{\infty} a_n$ and $\sum_{n=0}^{\infty} b_n$ be convergent series. Then:

    \[
        \sum_{n=0}^{\infty} (a_n + b_n)
        =
        \sum_{n=0}^{\infty} a_n
        +
        \sum_{n=0}^{\infty} b_n, \quad
        \sum_{n=0}^{\infty} C a_n
        =
        C \sum_{n=0}^{\infty} a_n
        \quad (C \in \mathbb{R}).
    \]
}
\\
\tsLem{(Tail of a Series)}{
    Let $\sum_{n=0}^{\infty} a_n$ be a series. For any $N \in \mathbb{N}_0$:
    \(
        \sum_{n=0}^{\infty} a_n \text{ converges}
        \quad \Leftrightarrow \quad
        \sum_{n=N}^{\infty} a_n \text{ converges},
    \)
    and in that case
    \[
        \sum_{n=0}^{\infty} a_n
        =
        \sum_{n=0}^{N-1} a_n
        +
        \sum_{n=N}^{\infty} a_n.
    \]
}
\\
\tsThe{(Series with Non-Negative Terms)}{
    Let $a_n \ge 0$ for all $n$. Then the sequence of partial sums
    \(
        s_n = \sum_{k=0}^{n} a_k
    \)
    is monotonically increasing. If $(s_n)$ is bounded, then the series $\sum_{n=0}^{\infty} a_n$ converges; otherwise, it diverges.
}
\\
\tsThe{(Comparison Test)}{
    Let $\sum a_n$, $\sum b_n$, $\sum c_n$ be series with non-negative terms and suppose that for some $N$:
    \(
        c_n \le b_n \le a_n \quad \text{for all } n \ge N.
    \)
    Then:
    \begin{itemize}
        \item If $\sum a_n$ converges, then $\sum b_n$ converges (majorant test).
        \item If $\sum c_n$ diverges, then $\sum b_n$ diverges (minorant test).
    \end{itemize}
}
\tsDef{(Absolute and Conditional Convergence)}{
    A series $\sum a_n$ is called \emph{absolutely convergent} if
    \(
    \sum |a_n| \text{ converges}.
    \)
    $\tsPoint$ It is called \emph{conditionally convergent} if $\sum a_n$ converges but $\sum |a_n|$ does not.
}
\\
\tsThe{(Riemann Rearrangement Theorem)}{
    Let $\sum a_n$ be a conditionally convergent series of real numbers, $L \in \mathbb{R}$. 
    Then there exists a bijection $\varphi : \mathbb{N}_0 \to \mathbb{N}_0$ such that
    \(
        \sum_{n=0}^{\infty} a_{\varphi(n)} = L.
    \)
}
\tsThe{(Rearrangement of Absolutely Convergent Series)}{
    Let $\sum a_n$ be absolutely convergent and let $\varphi : \mathbb{N}_0 \to \mathbb{N}_0$ be a bijection. Then
    \(
        \sum_{n=0}^{\infty} a_{\varphi(n)}
    \)
    converges absolutely and
    \(
        \sum_{n=0}^{\infty} a_n
        =
        \sum_{n=0}^{\infty} a_{\varphi(n)}.
    \)
}
\tsThe{(Leibniz Criterion)}{
    Let $(a_n)$ be a monotonically decreasing sequence of non-negative real numbers with
    \[
        \lim_{n\to\infty} a_n = 0.
    \]
    Then the alternating series
    \[
        \sum_{n=0}^{\infty} (-1)^n a_n
    \]
    converges. Moreover, for all $m \in \mathbb{N}_0$:
    \[
        \sum_{n=0}^{2m+1} (-1)^n a_n
        \le
        \sum_{n=0}^{\infty} (-1)^n a_n
        \le
        \sum_{n=0}^{2m} (-1)^n a_n.
    \]
}
\tsThe{(Cauchy Criterion for Series)}{
    The series $\sum_{n=0}^{\infty} a_n$ converges if and only if for every $\varepsilon > 0$ there exists $N$ such that for all $n \ge m \ge N$:
    \[
        \left| \sum_{k=m+1}^{n} a_k \right| \le \varepsilon.
    \]
}
\tsThe{(Absolute Convergence Implies Convergence)}{
    Every absolutely convergent series converges, and
    \[
        \left| \sum_{n=0}^{\infty} a_n \right|
        \le
        \sum_{n=0}^{\infty} |a_n|.
    \]
}
\tsThe{(Root Test)}{
    Let $(a_n)$ be a sequence and define
    \(
        \rho = \limsup_{n\to\infty} |a_n|^{1/n}.
    \)
    \\
    Then:
    \begin{itemize}[noitemsep]
        \item If $\rho < 1$, the series $\sum a_n$ converges absolutely.
        \item If $\rho > 1$, the series $\sum a_n$ diverges.
    \end{itemize}
}
\tsThe{(Ratio Test)}{
    Let $(a_n)$ with $a_n \neq 0$ and define
    \(
        \rho = \lim_{n\to\infty} \left| \frac{a_{n+1}}{a_n} \right|.
    \)
    Then:
    \begin{itemize}[noitemsep]
        \item If $\rho < 1$, the series $\sum a_n$ converges absolutely.
        \item If $\rho > 1$, the series $\sum a_n$ diverges.
    \end{itemize}
}
\tsThe{(Cauchy Product)}{
    Let $\sum a_n$ and $\sum b_n$ be absolutely convergent. Then
    \[
        \left( \sum_{n=0}^{\infty} a_n \right)
        \left( \sum_{n=0}^{\infty} b_n \right)
        =
        \sum_{n=0}^{\infty} \left( \sum_{k=0}^{n} a_{n-k} b_k \right),
    \]
    and the series on the right converges absolutely.
}
\\
\tsDef{(Power Series)}{
    A series of the form
    \(
        \sum_{k=0}^{\infty} c_k (x-a)^k
    \)
    is called a \emph{power series} centered at $a$ with coefficients $c_k$.
}
\\
\tsDef{(Convergence of a Power Series)}{
    For fixed $x$, if the sequence of partial sums
    \(
        s_n(x) = \sum_{k=0}^{n} c_k (x-a)^k
    \)
    converges, then the power series is said to converge at $x$.
}
\\
\tsThe{(Radius of Convergence)}{
    Every power series has a radius of convergence $R$ such that:
    \begin{itemize}[noitemsep]
        \item The series converges for $|x-a| < R$,
        \item The series diverges for $|x-a| > R$.
    \end{itemize}
    The interval $(a-R, a+R)$ is called the interval of convergence. The radius $R$ is given by
    \(
        \rho = \limsup_{n\to\infty} |c_n|^{1/n}
    \)
    , then:
    \[
        R =
        \begin{cases}
            0              & \text{if } \rho = \infty,     \\
            \frac{1}{\rho} & \text{if } 0 < \rho < \infty, \\
            \infty         & \text{if } \rho = 0.
        \end{cases}
    \]
}